\chapter{Introduction}\label{chap:introduction}

\section{Context and Motivation}\label{sec:intro-context}

The rapid expansion of DNA sequencing technologies has increased the demand for
computational methods capable of transforming raw genomic sequences into
biologically meaningful annotations. Functional genome annotation includes the
identification of DNA regions associated with regulatory processes, including
promoter regions. Promoters are located near the beginning of genes and are
involved in transcription initiation, making their identification relevant for
understanding gene regulation.

Promoter prediction can be formulated as a binary classification problem in
which a DNA sequence window is assigned to either the promoter or non-promoter
class. This thesis studies this task using machine learning and deep learning
models, with a particular focus on capsule-based neural architectures and
feature representations derived from nucleotide composition and physicochemical
properties.

\section{Research Problem}\label{sec:intro-problem}

Despite several promoter prediction methods reported in the literature,
accurately identifying promoters across heterogeneous organisms and data sources
remains challenging. Datasets differ in sequence length, transcription start
site annotation, organism group, negative-sample construction, and class balance.
These differences motivate a reproducible experimental protocol in which preprocessing, feature construction, model comparison, and reporting are specified explicitly.

\section{Hypothesis}\label{sec:intro-hypothesis}

The central hypothesis of this thesis is that capsule-based neural networks,
when evaluated under a reproducible experimental protocol and combined with
appropriate DNA representations, can capture structured sequence patterns that
are useful for promoter region prediction.

\section{Objectives}\label{sec:intro-objectives}

The general objective is to develop, evaluate, and document reproducible machine
learning approaches for promoter prediction in DNA sequences, with emphasis on
CapsProm-derived architectures.

The specific objectives are:

\begin{enumerate}
    \item Define a reproducible preprocessing and validation protocol for the active promoter datasets.
    \item Evaluate multiple DNA sequence representations, including one-hot, k-mer, and physicochemical encodings.
    \item Reproduce and organize baseline experiments under a unified protocol.
    \item Develop and evaluate capsule-based architectures for promoter prediction.
    \item Investigate extensions based on attention mechanisms, uncertainty estimation, and interpretability.
\end{enumerate}

\section{Thesis Organization}\label{sec:intro-organization}

Chapter~\ref{chap:background} introduces genomic and machine learning concepts
needed for the thesis. Chapter~\ref{chap:related-work} reviews promoter
prediction methods and prior work that motivates this research. Chapter~\ref{chap:materials-methods}
describes the datasets, preprocessing workflow, and reproducible experimental protocol.
Chapter~\ref{chap:feature-engineering} presents the feature engineering
strategies. Chapter~\ref{chap:models} describes baseline and proposed models.
Chapter~\ref{chap:experimental-protocol} defines the evaluation methodology.
Chapter~\ref{chap:results} reports the experimental results. Chapter~\ref{chap:discussion}
discusses the findings and limitations. Chapter~\ref{chap:conclusion} concludes
the thesis and outlines future work.
